\documentclass[11pt]{letter}
\usepackage{geometry}
\geometry{letterpaper, margin=1in}

\address{
[Author Name]\\
[Author Affiliation]\\
[Author Address]\\
[Author Email]
}

\begin{document}

\begin{letter}{
Editor-in-Chief\\
Computer Methods in Applied Mechanics and Engineering (CMAME)\\
Elsevier
}

\opening{Dear Editor-in-Chief,}

I am writing to submit our manuscript entitled \textbf{"Zugzwang: The Structural Anatomy of Physics-Informed Neural Network Failure"} for consideration for publication in \textit{Computer Methods in Applied Mechanics and Engineering} (CMAME).

In this paper, we address a critical open problem in scientific machine learning: why Physics-Informed Neural Networks (PINNs) reliably fail on stiff, advection-dominated partial differential equations despite extensive hyperparameter tuning. While recent literature typically attributes these failures to isolated phenomena such as gradient pathologies or precision errors, we introduce the \textit{Zugzwang} hypothesis to demonstrate that these pathologies are deeply interconnected structural traps.

Through a massive empirical framework comprising 24 controlled experiments across the 1D Advection, Viscous Burgers, and 2D Helmholtz equations, we quantitatively map out an "Atlas" of PINN failure. Our central practical finding directly addresses the CMAME applied-mechanics readership: \textbf{simply scaling up collocation points or internal network capacity cannot rescue PINNs from stiffness-induced failure}. By treating PINN convergence as a thermodynamic phase state, we demonstrate that the transition from success to catastrophic failure is a steep convergence cliff, fundamentally driven by the initial eigenvalue spectrum of the Neural Tangent Kernel.

We believe our rigorous experimental approach, the unified phase-space framework, and the explicit variance decomposition analysis identifying the structural dominance of the spectral barrier will be of high interest to the readership of CMAME, particularly those developing robust numerical-ML hybrid methods for complex continuum mechanics applications.

This manuscript is original, has not been published previously, and is not under consideration for publication elsewhere. All data and code required to reproduce the findings are fully documented and publicly available.

Thank you for your consideration of this work.

\closing{Sincerely,\\
\vspace{1cm}
[Author Name]
}

\end{letter}
\end{document}
